\clearpage
\SourceLeaf{264}
\thispagestyle{empty}
\pdfbookmark[0]{Investigations into Topics in Higher Geodesy. First Memoir}{geodaesie1}
\vspace*{.3in}
\begin{center}
{\large INVESTIGATIONS}\par\bigskip
INTO\par\bigskip
{\Large TOPICS IN HIGHER GEODESY}\par\vspace{25pt}
FIRST MEMOIR\par\vspace{35pt}
BY\par\vspace{35pt}
{\large CARL FRIEDRICH GAUSS}\par\vspace{35pt}
{\small PRESENTED TO THE ROYAL SOCIETY ON 23 OCTOBER 1843.}
\end{center}
\vfill
\begin{center}
\rule{.8\linewidth}{.3pt}\par\smallskip
Abhandlungen der Königl. Gesellschaft der Wissenschaften zu Göttingen. Vol. II.\par
Göttingen 1844.\par\smallskip
\rule{.8\linewidth}{.3pt}
\end{center}
\par\bigskip\begingroup\fontsize{9}{11}\selectfont\raggedright\noindent The first higher-geodesy memoir is translated from the German in \textit{Werke}, Band IV. Its appended tables are reproduced from the same edition. The opening of the \href{https://doi.org/10.3931/e-rara-2860}{ETH-Bibliothek copy of the separate 1844 printing} was consulted in OCR. Related terminology was compared with \href{https://arxiv.org/abs/math/0409578}{Erhard Scholz} (2004), \href{https://doi.org/10.12902/zfv-0387-2022}{Miljenko Lapaine} (2022), and \href{https://en.wikisource.org/w/index.php?oldid=12858814}{J. W. L. Glaisher's entry on mathematical tables} (1911).\par\endgroup
\clearpage
